% ================================================================
\section{The Mertens Function as Magnetization}
\label{sec:mertens}
% ================================================================

The Mertens function $M(x) = \sum_{n \leq x} \mu(n)$ plays a central
role in the forward direction of the proof.

\begin{correspondence}[Ising Model]
The M\"obius function $\mu(n) \in \{-1, 0, +1\}$ behaves like a
\emph{spin variable} in the Ising model:
\begin{itemize}
  \item $\mu(n) = +1$: spin up (squarefree, even number of prime factors)
  \item $\mu(n) = -1$: spin down (squarefree, odd number of prime factors)
  \item $\mu(n) = 0$: vacancy (contains a squared prime factor)
\end{itemize}
The Mertens function $M(x) = \sum_{n \leq x} \mu(n)$ is the net
\emph{magnetization} up to $x$.

The Riemann Hypothesis is equivalent to $M(x) = O(x^{1/2+\varepsilon})$---the
magnetization fluctuates like $\sqrt{x}$, which is the \emph{central limit
theorem} behavior expected for a system with short-range correlations at
its critical point.
\end{correspondence}

The Cathedral axiom \lean{rh\_implies\_mertens\_bound} states:
\[
  \RH \implies |M(x)| \leq C \cdot x^{1/2} \log^2 x
\]
This is the physics statement that \emph{under RH, the prime spin system
is at its critical temperature}---the fluctuations are exactly
$\Theta(\sqrt{x})$, neither ordered (like $O(1)$, which would give PNT
with power-saving error) nor disordered (like $O(x^{3/4})$, which would
violate RH).

The weaker bound $|M(x)| \leq 64\, C \cdot x^{3/4}$ is a \emph{proved
corollary} (\lean{rh\_implies\_mertens\_34}, April 22, 2026)---the Mertens
magnetization is bounded in any sub-critical regime.

\subsection{Abel Summation as Renormalization}

The Abel summation step in the forward proof---converting the Mertens
bound into an $L^2$ bound on the M\"obius weights---is the
\emph{renormalization} procedure:

\[
  v_k = -\mu(k)\left(1 - \frac{\ln k}{\ln N}\right)
\]

The log-linear taper $1 - \ln k / \ln N$ is a \emph{UV cutoff}: it
smoothly suppresses the contribution of high-$k$ modes (analogous to
high-momentum states in QFT). Without this cutoff, the weight norm
$\|v\|^2 = \Theta(N)$ diverges---a UV divergence that the taper
renormalizes away.

\begin{observation}[The Triangle Inequality Trap as Failed Renormalization]
The Cathedral's Discovery~5 (the Triangle Inequality Trap) identifies
a situation where naive bounds destroy the cancellation needed for
$d_N^2 \to 0$. In physics language: applying the triangle inequality
to $E = 1 - 2b^Tv + v^TGv$ is like computing the energy without
accounting for interference---it gives $E \geq 4$ for a quantity
approaching $0$. This is precisely the problem that renormalization
was invented to solve: divergent intermediate quantities that cancel
in the final answer.

The Cathedral resolves this trap in \lean{MoebiusL1Bound.lean} by
decomposing $1 - b^T v$ into the thermodynamic hierarchy (\S\ref{sec:thermo}),
avoiding the triangle inequality entirely.
\end{observation}

\subsection{The Thermodynamic Hierarchy}
\label{sec:thermo}

The proved identity (\lean{CalcBounds.lean}, zero sorry):
\[
  1 - b^T v
  = (1 - \gamma)\, S_1 + (S_2 + 1)
  - \frac{(1 - \gamma)\,S_2 + S_3}{\ln N}
\]
where $S_m = \sum_{k=1}^{N-1} \mu(k)(\ln k)^{m-1}/k$ are the Abel
partial sums, is a \emph{moment expansion} of the grand partition function
$1/\zeta(s)$ near $s = 1$.

\begin{correspondence}[Thermodynamic Moments]
  Each Abel sum $S_m$ corresponds to a response function of the prime
  spin system:
  \begin{center}
  \begin{tabular}{@{}lll@{}}
  \toprule
  \textbf{Abel Sum} & \textbf{Thermodynamic Dual} & \textbf{Limiting Value} \\
  \midrule
  $S_1 = \sum \mu(k)/k$ & Magnetization (PNT) & $\to 0$ \\
  $S_2 = \sum \mu(k)\ln k / k$ & Magnetic susceptibility & $\to -1$ \\
  $S_3 = \sum \mu(k)\ln^2 k / k$ & Heat capacity & $\to -2\gamma$ \\
  \bottomrule
  \end{tabular}
  \end{center}
  The correction term $[(1-\gamma)S_2 + S_3]/\ln N$ is the
  \emph{finite-size scaling} contribution---the deviation from the
  thermodynamic limit due to the cutoff at $N$.
\end{correspondence}

The Cathedral proves each $S_m$ decays with explicit rates
(\lean{AbelTail/*.lean}, zero sorry):
$|S_1| \leq C_1 \cdot N^{-1/4}$,
$|S_2 + 1| \leq C_2 \cdot N^{-1/4} \ln N$,
$|S_3 + 2\gamma|$ bounded universally.
These are the \emph{critical exponents} of the prime number gas.
Numerical validation (256-bit MPFR, \texttt{millennium-wall} experiment)
confirms $|1 - b^T v| \cdot \ln N \to \gamma + 1 \approx 1.577$---the
\emph{Euler--Mascheroni amplitude} governs the rate of vacuum screening.

